%! Author = sriadityavedantam
%! Date = 3/30/26

\section{Galois Theory}

Define a homomorphism
\[
    \phi:\mathbb{Z}\to \mathbb{F}
\]
by
\[
    1\mapsto 1
\]
and hence
\[
    n\mapsto n1.
\]
Then
\[
    \phi(n+m)=\phi(n)+\phi(m)
\]
and
\[
    \phi(nm)=\phi(n)\phi(m).
\]
Thus $\mathrm{Im}(\phi)$ is a subring of $\mathbb{F}$.
Also,
\[
    \frac{\mathbb{Z}}{\ker\phi}\cong \mathrm{Im}(\phi).
\]
Since $\ker\phi$ is an ideal of $\mathbb{Z}$, it has the form
\[
    n\mathbb{Z}
\]
for some $n\geq 0$.
If
\[
    \ker\phi=\{0\},
\]
then $\mathbb{F}$ contains a subring isomorphic to $\mathbb{Z}$, hence a subfield isomorphic to $\mathbb{Q}$.
We say in this case that $\mathbb{F}$ has characteristic $0$.
If
\[
    \ker\phi=p\mathbb{Z}
\]
for some prime $p$, then
\[
    \mathbb{F}
\]
contains a subfield isomorphic to
\[
    \mathbb{Z}_p\cong \frac{\mathbb{Z}}{p\mathbb{Z}}.
\]
We then say that $\mathbb{F}$ has characteristic $p$.

\begin{definition}[Formal Derivative]
    If
    \[
        f(x)\in \mathbb{F}[x],
    \]
    then the formal derivative of $f$ is denoted by
    \[
        f'(x)=Df(x).
    \]
    If
    \[
        f(x)=x^n,
    \]
    then
    \[
        Df(x)=nx^{n-1},
    \]
    and this extends to all of $\mathbb{F}[x]$ by linearity.
\end{definition}

Alternatively, in $\mathbb{F}[x,h]$, where $h$ is transcendental, define
\[
    g(x,h)=f(x+h)-f(x).
\]
Then
\[
    g(x,0)=0.
\]
By the factor theorem, $h$ divides $g(x,h)$.

\begin{definition}
    \[
        Df(x)
    \]
    is the polynomial obtained from
    \[
        \frac{g(x,h)}{h}
    \]
    by evaluating at
    \[
        h=0.
    \]
\end{definition}

\begin{definition}[Repeated Root]
    Let
    \[
        f(x)\in \mathbb{F}[x],
    \]
    and let $\mathbb{E}$ be a splitting field for $f(x)$.
    We say that $a$ is a repeated root of $f(x)$ if
    \[
        f(x)=(x-a)^{2}g(x)
    \]
    in $\mathbb{E}[x]$ for some $g(x)\in \mathbb{E}[x]$.
\end{definition}

If $a$ is a repeated root of $f(x)$, then
\[
    f(a)=0
    \quad\text{and}\quad
    Df(a)=0.
\]
So $a$ is a common root of $f(x)$ and $Df(x)$.
Now assume $\mathbb{E}$ is a finite-dimensional extension of $\mathbb{Q}$.
Then $\mathrm{char}(\mathbb{E})=0$.
Let $\mathbb{F}$ be a subfield of $\mathbb{E}$.
If
\[
    f(x)\in \mathbb{F}[x]
\]
is irreducible and $u\in \mathbb{E}$ is a root of $f(x)$, then $u$ is not a repeated root of $f(x)$.
Indeed, if $u$ were a repeated root, then
\[
    Df(u)=0.
\]
Since
\[
    \deg(Df)<\deg(f),
\]
this would contradict the minimality of the irreducible polynomial $f(x)$.

\begin{definition}[Normal Extension]
    Suppose $\mathbb{E}/\mathbb{F}$ is a field extension.
    We say $\mathbb{E}$ is normal over $\mathbb{F}$ if whenever
    \[
        u\in \mathbb{E}
    \]
    and $f(x)\in \mathbb{F}[x]$ is the minimal polynomial of $u$ over $\mathbb{F}$, then $f(x)$ splits completely in $\mathbb{E}[x]$.
\end{definition}

Equivalently, any polynomial in $\mathbb{F}[x]$ which has one root in $\mathbb{E}$ has all of its roots in $\mathbb{E}$.

\begin{example}
    If $\mathbb{E}$ is a splitting field for
    \[
        f(x)\in \mathbb{F}[x],
    \]
    then $\mathbb{E}$ is normal over $\mathbb{F}$.
    For example, let
    \[
        \omega=e^{2\pi i/3}.
    \]
    Then
    \[
        \mathbb{Q}\left(2^{1/3},\,2^{1/3}\omega,\,2^{1/3}\omega^2\right)
        =
        \mathbb{Q}\left(2^{1/3},\omega\right)
    \]
    is the splitting field for
    \[
        x^3-2
    \]
    over $\mathbb{Q}$, so it is normal over $\mathbb{Q}$.
    But
    \[
        \mathbb{Q}\left(2^{1/3}\right)
    \]
    is not normal over $\mathbb{Q}$.
\end{example}

\begin{definition}[Automorphism]
    \[
        \aut(\mathbb{E})
    \]
    is the group of field automorphisms of $\mathbb{E}$.
\end{definition}

\begin{definition}[Galois Group]
    \[
        \aut\left(\frac{\mathbb{E}}{\mathbb{F}}\right)
        \coloneqq
        \{\sigma\in \aut(\mathbb{E}):\sigma(a)=a\ \forall a\in \mathbb{F}\}
        =
        \gal\left(\frac{\mathbb{E}}{\mathbb{F}}\right).
    \]
\end{definition}

\begin{definition}[Fixed Field]
    Let
    \[
        G=\aut\left(\frac{\mathbb{E}}{\mathbb{F}}\right).
    \]
    The fixed field of $G$ is
    \[
        \{u\in \mathbb{E}:\sigma(u)=u\ \forall \sigma\in G\}
        =
        \inv(G).
    \]
\end{definition}

Note that
\[
    \mathbb{F}\subseteq \inv(G),
\]
but it is not automatic that $\inv(G)=\mathbb{F}$.

\begin{example}
    Let
    \[
        \mathbb{F}=\mathbb{Q},
        \qquad
        \mathbb{E}=\mathbb{Q}(2^{1/3}).
    \]
    Then
    \[
        G=\aut\left(\frac{\mathbb{E}}{\mathbb{F}}\right)=\{1\}.
    \]
    So
    \[
        \inv(G)=\mathbb{E},
    \]
    which is strictly larger than $\mathbb{F}$.
\end{example}

\begin{theorem}[Fundamental Theorem of Galois Theory]{
    Assume
    \[
        [\mathbb{E}:\mathbb{F}]<\infty
    \]
    and
    \[
        \mathrm{char}(\mathbb{F})=0.
    \]
    Then the following are equivalent:
    \begin{enumerate}
        \item
        \[
            \inv\left(\aut\left(\frac{\mathbb{E}}{\mathbb{F}}\right)\right)=\mathbb{F};
        \]
        \item
        \[
            \mathbb{F}=\inv(G)
        \]
        for some finite subgroup
        \[
            G\leq \aut(\mathbb{E});
        \]
        \item $\mathbb{E}$ is normal over $\mathbb{F}$;
        \item $\mathbb{E}$ is the splitting field of a separable polynomial in $\mathbb{F}[x]$;
        \item
        \[
            \ord\left(\aut\left(\frac{\mathbb{E}}{\mathbb{F}}\right)\right)
            =
            [\mathbb{E}:\mathbb{F}].
        \]
    \end{enumerate}
}
    We sketch the main implications.
    \[
        \mathbf{(4)}\implies \mathbf{(5)}
    \]
    has already been shown.
    \[
        \mathbf{(1)}\implies \mathbf{(2)}
    \]
    is immediate by taking
    \begin{gather*}
        G=\aut\left(\frac{\mathbb{E}}{\mathbb{F}}\right).\\
        \mathbf{(3)}\implies \mathbf{(4)}.\\
    \end{gather*}
    Since
    \[
        [\mathbb{E}:\mathbb{F}]<\infty,
    \]
    we can write
    \[
        \mathbb{E}=\mathbb{F}(u_1,\ldots,u_n).
    \]
    Let
    \[
        p_i(x)\in \mathbb{F}[x]
    \]
    be the minimal polynomial of $u_i$ over $\mathbb{F}$.
    Because $\mathbb{E}/\mathbb{F}$ is normal, each $p_i(x)$ splits in $\mathbb{E}[x]$.
    Let $f(x)$ be the product of the distinct $p_i(x)$.
    Then $\mathbb{E}$ is the splitting field of $f(x)$ over $\mathbb{F}$.
    Since $\mathrm{char}(\mathbb{F})=0$, the polynomial is separable.
    \[
        \mathbf{(2)}\implies \mathbf{(3)}.
    \]
    Assume
    \[
        \mathbb{F}=\inv(G)
    \]
    for some finite subgroup
    \[
        G\leq \aut(\mathbb{E}).
    \]
    Let $u\in \mathbb{E}$.
    Consider its orbit under $G$:
    \[
        u_1=u,\ u_2,\ldots,u_m.
    \]
    Define
    \[
        g(x)=(x-u_1)\cdots(x-u_m)\in \mathbb{E}[x].
    \]
    We claim that
    \[
        g(x)\in \mathbb{F}[x].
    \]
    Indeed, each $\sigma\in G$ permutes the roots $u_i$, so it fixes the coefficients of $g(x)$.
    Hence the coefficients lie in $\inv(G)=\mathbb{F}$.
    Thus the minimal polynomial of $u$ over $\mathbb{F}$ divides $g(x)$ and therefore splits in $\mathbb{E}[x]$.
    So $\mathbb{E}/\mathbb{F}$ is normal.
    Since $\mathrm{char}(\mathbb{F})=0$, it is also separable.
    \[
        \mathbf{(5)}\implies \mathbf{(1)}.
    \]
    Let
    \[
        \mathbb{K}=\inv\left(\aut\left(\frac{\mathbb{E}}{\mathbb{F}}\right)\right).
    \]
    Then by construction
    \[
        \mathbb{F}\subseteq \mathbb{K}\subseteq \mathbb{E}.
    \]
    Also,
    \[
        \aut\left(\frac{\mathbb{E}}{\mathbb{K}}\right)
        =
        \aut\left(\frac{\mathbb{E}}{\mathbb{F}}\right).
    \]
    By Artin's theorem and condition \(\mathbf{(5)}\),
    \[
        [\mathbb{E}:\mathbb{K}]
        =
        \ord\left(\aut\left(\frac{\mathbb{E}}{\mathbb{K}}\right)\right)
        =
        \ord\left(\aut\left(\frac{\mathbb{E}}{\mathbb{F}}\right)\right)
        =
        [\mathbb{E}:\mathbb{F}].
    \]
    Hence
    \[
        \mathbb{K}=\mathbb{F}.
    \]
    So condition \(\mathbf{(1)}\) holds.
\end{theorem}

\begin{lemma}{
    Let $\mathbb{E}$ be an extension field of $\mathbb{F}$.
    Let
    \[
        G=\aut\left(\frac{\mathbb{E}}{\mathbb{F}}\right)
    \]
    and let
    \[
        \mathbb{K}=\inv(G).
    \]
    Then
    \[
        \aut\left(\frac{\mathbb{E}}{\mathbb{K}}\right)
        =
        \aut\left(\frac{\mathbb{E}}{\mathbb{F}}\right)
        =
        G.
    \]
}
    This is immediate from the definition of fixed field.
    An automorphism fixes $\mathbb{F}$ if and only if it fixes every element fixed by all elements of $G$.
\end{lemma}

\begin{definition}[Galois]
    We say
    \[
        \frac{\mathbb{E}}{\mathbb{F}}
    \]
    is Galois if any of the equivalent conditions above hold.
\end{definition}

\begin{theorem}{
    If
    \[
        G\leq \aut(\mathbb{E})
    \]
    is finite and
    \[
        \mathbb{K}=\inv(G),
    \]
    then
    \[
        [\mathbb{E}:\mathbb{K}]\leq \ord(G).
    \]
}
    This is Artin's Lemma.
    Suppose, for contradiction, that there exist
    \[
        u_1,\ldots,u_m\in \mathbb{E}
    \]
    which are linearly independent over $\mathbb{K}$ with
    \[
        m>\ord(G)=n.
    \]
    Write
    \[
        G=\{\sigma_1,\ldots,\sigma_n\},
    \]
    where $\sigma_1=1$.
    Form the matrix
    \[
        A=
        \begin{bmatrix}
            \sigma_1(u_1) & \sigma_1(u_2) & \cdots & \sigma_1(u_m)\\
            \sigma_2(u_1) & \sigma_2(u_2) & \cdots & \sigma_2(u_m)\\
            \vdots & \vdots & & \vdots\\
            \sigma_n(u_1) & \sigma_n(u_2) & \cdots & \sigma_n(u_m)
        \end{bmatrix}.
    \]
    This is an $n\times m$ matrix with $m>n$, so the system
    \[
        A\vec{x}=\vec{0}
    \]
    has a nontrivial solution.
    Choose a nontrivial solution
    \[
        \vec{b}=
        \begin{bmatrix}
            b_1\\
            \vdots\\
            b_m
        \end{bmatrix}
    \]
    with the fewest number of nonzero entries.
    After scaling and reordering, we may assume
    \[
        b_1=1.
    \]
    We claim that each
    \[
        b_i\in \mathbb{K}.
    \]
    Suppose not.
    Then there exists some $\sigma_\ell\in G$ such that
    \[
        \sigma_\ell(b_j)\neq b_j
    \]
    for some $j$.
    Applying $\sigma_\ell$ to the equation
    \[
        A\vec{b}=\vec{0}
    \]
    shows that
    \[
        \sigma_\ell(\vec{b})
    \]
    is also a solution.
    Hence
    \[
        \vec{b}-\sigma_\ell(\vec{b})
    \]
    is a nontrivial solution with fewer nonzero entries, contradicting minimality.
    Thus
    \[
        \vec{b}\in \mathbb{K}^m.
    \]
    So we have found a nontrivial linear relation among
    \[
        u_1,\ldots,u_m
    \]
    with coefficients in $\mathbb{K}$, contradicting linear independence.
    Therefore
    \[
        [\mathbb{E}:\mathbb{K}]\leq \ord(G).
    \]
\end{theorem}

\begin{theorem}{
    Suppose
    \[
        \mathbb{E}/\mathbb{F}
    \]
    is Galois, and let
    \[
        G=\aut\left(\frac{\mathbb{E}}{\mathbb{F}}\right).
    \]
    Then for every intermediate field
    \[
        \mathbb{F}\subseteq \mathbb{K}\subseteq \mathbb{E},
    \]
    the extension
    \[
        \mathbb{E}/\mathbb{K}
    \]
    is Galois.
    Moreover, there is a bijection between intermediate fields $\mathbb{K}$ and subgroups
    \[
        H\leq G.
    \]
    The correspondence is given by
    \[
        \mathbb{K}\longmapsto \aut\left(\frac{\mathbb{E}}{\mathbb{K}}\right)
    \]
    and
    \[
        H\longmapsto \inv(H).
    \]
    These maps are inclusion-reversing.
}
    Since $\mathbb{E}/\mathbb{F}$ is Galois, it is the splitting field of some separable polynomial
    \[
        f(x)\in \mathbb{F}[x].
    \]
    Because
    \[
        \mathbb{F}\subseteq \mathbb{K},
    \]
    the same polynomial lies in $\mathbb{K}[x]$, and $\mathbb{E}$ is still its splitting field.
    Thus
    \[
        \mathbb{E}/\mathbb{K}
    \]
    is Galois.
    Now let
    \[
        H\leq G
    \]
    and define
    \[
        \mathbb{K}=\inv(H).
    \]
    By Artin's Lemma,
    \[
        [\mathbb{E}:\mathbb{K}]\leq \ord(H).
    \]
    Also,
    \[
        H\leq \aut\left(\frac{\mathbb{E}}{\mathbb{K}}\right).
    \]
    Since
    \[
        \mathbb{E}/\mathbb{K}
    \]
    is Galois, we have
    \[
        \ord\left(\aut\left(\frac{\mathbb{E}}{\mathbb{K}}\right)\right)
        =
        [\mathbb{E}:\mathbb{K}].
    \]
    So
    \[
        [\mathbb{E}:\mathbb{K}]
        \leq
        \ord(H)
        \leq
        \ord\left(\aut\left(\frac{\mathbb{E}}{\mathbb{K}}\right)\right)
        =
        [\mathbb{E}:\mathbb{K}].
    \]
    Hence
    \[
        H=\aut\left(\frac{\mathbb{E}}{\mathbb{K}}\right).
    \]
    Thus the two maps are inverses.
\end{theorem}

These maps are order-reversing under inclusion.
That is,
\[
    \mathbb{K}_1\subseteq \mathbb{K}_2
    \implies
    \aut\left(\frac{\mathbb{E}}{\mathbb{K}_1}\right)\supseteq
    \aut\left(\frac{\mathbb{E}}{\mathbb{K}_2}\right),
\]
and if
\[
    H_1\subseteq H_2,
\]
then
\[
    \inv(H_1)\supseteq \inv(H_2).
\]

\begin{example}[Galois Groups Fork]
    Let
    \[
        \mathbb{F}=\mathbb{Q}
    \]
    and
    \[
        \mathbb{E}=\mathbb{Q}\left(2^{1/3},\omega\right),
        \qquad
        \omega=e^{2\pi i/3}.
    \]
    Then $\mathbb{E}$ is Galois over $\mathbb{Q}$ because it is the splitting field of
    \[
        f(x)=x^3-2
    \]
    over $\mathbb{Q}$.
    Also,
    \[
        [\mathbb{Q}(\omega):\mathbb{Q}]=2.
    \]
    When $\mathbb{E}/\mathbb{F}$ is Galois, the intermediate fields and the subgroups of
    \[
        G=\aut\left(\frac{\mathbb{E}}{\mathbb{F}}\right)
    \]
    correspond exactly:
    \begin{gather*}
        \mathbb{E}
        \stackrel{\aut(\mathbb{E}/\mathbb{E})}{\underset{\mathbb{E}=\inv(1)}{\lrarrow}}
        1,\\
        \mathbb{K}
        \stackrel{\aut(\mathbb{E}/\mathbb{K})}{\underset{\mathbb{K}=\inv(H)}{\lrarrow}}
        H,\\
        \mathbb{F}
        \stackrel{\aut(\mathbb{E}/\mathbb{F})}{\underset{\mathbb{F}=\inv(G)}{\lrarrow}}
        G.\\
    \end{gather*}
    These maps are inverses.
\end{example}

What is the degree of
\[
    \mathbb{Q}\left(2^{1/3},\omega\right)?
\]
We get this field from the polynomial
\[
    x^3-2,
\]
which has $3$ roots.
Automorphisms permute the roots, so the Galois group has size at most
\[
    \ord(S_3)=6.
\]
Let
\[
    \mathbb{E}=\mathbb{Q}\left(2^{1/3},\omega\right).
\]
What are the automorphisms in
\[
    G=\gal\left(\frac{\mathbb{E}}{\mathbb{Q}}\right)?
\]
One such automorphism is complex conjugation.
Let
\[
    \sigma(z)=\overline{z}.
\]
Then
\[
    \overline{z+w}=\overline{z}+\overline{w}
\]
and
\[
    \overline{zw}=\overline{z}\,\overline{w}.
\]
Also,
\[
    \sigma\left(2^{1/3}\right)=2^{1/3}
\]
because
\[
    2^{1/3}\in \mathbb{R},
\]
and
\[
    \sigma(\omega)=\overline{\omega}=\omega^{-1}=\omega^2.
\]
Then
\[
    \sigma^2(\omega)=\sigma(\omega^2)=\omega,
\]
so
\[
    \sigma^2=1_G.
\]